\section{Бенчмарк на обычной нагрузке}

Исправление в \texttt{fixed}-версии касалось файлового ввода-вывода, поэтому отдельно было важно проверить, что типичная нагрузка из лабораторной работы \textnumero 2 не ухудшилась. Для этого обе версии были прогнаны через общий benchmark-workload на \(200000\) измеряемых операциях после предварительной вставки \(20000\) элементов. Для каждого сценария бралась медиана по пяти прогонам.

\begin{center}
\begin{tabular}{|l|r|r|r|}
\hline
Сценарий & Baseline, \texttt{us} & Fixed, \texttt{us} & Baseline / Fixed \\
\hline
\texttt{mixed} & \(50073\) & \(35497\) & \(1.41\times\) \\
\hline
\texttt{read-heavy} & \(48001\) & \(33159\) & \(1.45\times\) \\
\hline
\texttt{update-heavy} & \(39724\) & \(39384\) & \(1.01\times\) \\
\hline
\end{tabular}
\end{center}

Результаты показывают, что после исправления программа не только не замедлилась на обычной нагрузке, но и в двух сценариях заметно ускорилась. На \texttt{mixed} и \texttt{read-heavy} это выглядит как выигрыш порядка \(1.4\times\), а на \texttt{update-heavy} разница почти исчезает. Это соответствует ожиданию: исправление было локальным и не меняло алгоритмы работы самого дерева, но сняло часть лишних накладных расходов вокруг него.
\pagebreak
